% sec-02-introduction.tex - Introduction (draft scaffold)
\section{Introduction}

\subsection{Patients Currently Enrolled in Trial Cannot Wait for Slowdowns}
\draftinstr{Open with OUTLINE theme 1. Make the patient-centered case: enrolled
Phase 1 PDAC patients have under-13-percent five-year survival and cannot wait for
administrative slowdowns; faster paperwork enables faster treatment; new document
iterations now take 1-4 days; a single prompt grows a draft to a full to a final
document. Render \texttt{mermaid/fig-17-patient-time-saved-cascade.md} and
\texttt{fig-10-iteration-timeline.md} as TikZ figures. Source survival framing from
\texttt{inputs/references.bib} (Siegel-class statistics via the protocol works) and
the 1-4 day cadence from \texttt{prompts/prompt-paper.md} (OUTLINE).}

\subsection{Effective Verifications Prove Utility}
\draftinstr{Introduce OUTLINE theme 2: grounding via Mermaid figures; figures match
paper context; AI-generated files viewable on GitHub in real time for
monitorability; paper repository and directory names match the repository names;
paper file names match repo file names. Render
\texttt{mermaid/fig-15-name-matching-verification.md} and
\texttt{fig-14-grounding-mermaid-to-tikz.md}. Source from
\texttt{prompts/prompt-paper.md} (OUTLINE 2) and \texttt{inputs/llm-adoption/main.tex}
(monitorability).}

\subsection{Patient Lives Extended}
\draftinstr{Introduce OUTLINE theme 3: probable benefit greater than probable risk
(state exactly how), and the mechanism by which the method compresses
administrative/prep time. Render \texttt{mermaid/fig-16-benefit-risk-framework.md}
and \texttt{fig-11-time-bucket-compression.md}. Source the three timeline buckets
from \texttt{research/document-types/Gemini-3-1-Pro-DocTypes-2026-06-26.md}.}

\subsection{Scope, Contributions, and Relationship to the Adoption Guide}
\draftinstr{State that this paper builds on the Oncology Trial PI LLM Adoption Guide
\cite{adoptionguide} (the paper template, paper v1.0) and turns its hands-on
guidance into a worked Phase 1 PDAC document-generation study. List contributions:
(1) a single-prompt mermaid-draft-full-final pipeline; (2) coverage of the six
ACCELERATION targets; (3) 20+ grounded figures reproduced identically in LaTeX;
(4) a real-time monitorable build. Note the real-world daraxonrasib PDAC trial
thread \cite{phase1protocol, phase2protocol, kawchak_2026_20196639} and the
Physical AI framework \cite{natplatform}. State repository v4.2.0. Render
\texttt{mermaid/fig-01-single-prompt-pipeline.md} and
\texttt{fig-21-daraxonrasib-document-thread.md}.}
